%% Question bank chapter file - generated by bank_tools.py build (do not edit while the bank is still being built:
%% it is regenerated from the per-test data in Question Bank\data). Every question is one bankq block:
%%   \begin{bankq}[id=..., source={...}, paper=.., qno=.., type=..., subtopic={...}, answer={...}, answer_src=key|solved,
%%                 working={...}, difficulty=1-10, diff_parts={K S I C F}, diff_note={...},
%%                 group=<passage id or empty>, also_in={duplicate copies in other tests}]
%% difficulty = min(10, K+S+I+C+F-1): K concepts 1-3, S steps 1-4, I insight 0-2, C computation 0-2, F format 0-1
%%   ...question LaTeX, options, figures (\bfig / \bfigi -> img/)...
%%   \end{bankq}
%% Passages shared by several questions are bankpassage blocks; their questions carry group=<passage id>.
\documentclass[10pt]{article}
\usepackage[a4paper,margin=1.8cm]{geometry}
\usepackage[T1]{fontenc}
\usepackage{amsmath,amssymb,mathtools}
\usepackage{graphicx}
\usepackage[export]{adjustbox}
\usepackage{xcolor,needspace,keyval}
\usepackage{enumitem}
\usepackage{booktabs,array,tabularx}
\usepackage[hidelinks]{hyperref}
\setlength{\parindent}{0pt}
\setlength{\parskip}{3pt}
\makeatletter
\define@key{bq}{id}{\def\bq@id{#1}}\define@key{bq}{source}{\def\bq@src{#1}}\define@key{bq}{paper}{\def\bq@paper{#1}}
\define@key{bq}{qno}{\def\bq@qno{#1}}\define@key{bq}{type}{\def\bq@type{#1}}\define@key{bq}{subtopic}{\def\bq@sub{#1}}
\define@key{bq}{answer}{\def\bq@ans{#1}}\define@key{bq}{answer_src}{\def\bq@asrc{#1}}\define@key{bq}{working}{\def\bq@work{#1}}
\define@key{bq}{group}{\def\bq@group{#1}}\define@key{bq}{also_in}{\def\bq@also{#1}}
\define@key{bq}{difficulty}{\def\bq@diff{#1}}\define@key{bq}{diff_parts}{\def\bq@dparts{#1}}\define@key{bq}{diff_note}{\def\bq@dnote{#1}}
\newenvironment{bankq}[1][]{\def\bq@work{}\def\bq@also{}\def\bq@group{}\def\bq@diff{?}\def\bq@dparts{}\def\bq@dnote{}\setkeys{bq}{#1}%
  \par\needspace{4\baselineskip}\medskip\noindent{\small\textbf{\bq@id}\hfill\color{gray}[\bq@src, Paper \bq@paper, Q\bq@qno]
  \bq@type\ $\cdot$ \bq@sub\ $\cdot$ \textbf{D\bq@diff}}\par\nopagebreak\relax}
  {\par{\small\color{teal}Answer: \textbf{\bq@ans}\ifx\bq@work\@empty\else\ (solved: \bq@work)\fi
  \ifx\bq@also\@empty\else\par\color{gray}Also in: \bq@also\fi}\par}
\define@key{bp}{id}{\def\bp@id{#1}}\define@key{bp}{source}{\def\bp@src{#1}}\define@key{bp}{questions}{\def\bp@qs{#1}}
\newenvironment{bankpassage}[1][]{\setkeys{bp}{#1}\par\needspace{4\baselineskip}\bigskip\noindent
  {\small\textbf{Passage \bp@id}\hfill\color{gray}[\bp@src; questions \bp@qs]}\par\nopagebreak\relax\itshape}{\par}
\makeatother
\newcommand{\bfig}[2]{\par\smallskip{\centering\includegraphics[width=#1pt,max width=0.62\linewidth,max height=0.32\textheight]{img/#2}\par}\smallskip}
\newcommand{\bfigi}[2]{\includegraphics[width=#1pt,max width=\linewidth,max height=0.32\textheight,valign=c]{img/#2}}
\newcommand{\bmissing}[1]{\textit{[figure missing in the source conversion: \detokenize{#1}]}}
\newenvironment{qoptions}{\begin{itemize}[leftmargin=1.8em,itemsep=1pt,topsep=2pt,label={}]}{\end{itemize}}
\newcommand{\opt}[2]{\item[(#1)] #2}
\providecommand{\fwstart}{}\providecommand{\fwend}{}\providecommand{\qfullsplit}{}
\newcolumntype{L}[1]{>{\raggedright\arraybackslash}p{#1}}
\newcolumntype{C}[1]{>{\centering\arraybackslash}p{#1}}
\newcolumntype{R}[1]{>{\raggedleft\arraybackslash}p{#1}}
\providecommand{\sectionhdr}[1]{}\providecommand{\subsectionhdr}[1]{}


\begin{document}
{\Large\bfseries IUPAC --- Question Bank}\par\smallskip
{\small Chemistry chapter C26. Questions from all converted test papers, grouped by source series. Answers: printed key where it matches the paper, otherwise solved (answer\_src=solved, with working). Exact repeats across papers appear once, with also\_in listing the other copies.}\par\medskip\hrule

\section*{TOPS papers}
\begin{bankq}[id=TOPS03-P2-CHEM-S3-Q03, source={TOPS - Test 03 (20 Nov 2022)}, paper=2, qno=3, type=numerical, subtopic={Nomenclature rules}, answer={2}, answer_src=key, difficulty=3, diff_parts={K1 S1 I0 C1 F1}, diff_note={concepts: functional groups | steps: identify the groups in each molecule; count}]
Number of F.G. present in compound H$_2$N--CHO are $x$ and number of F.G. present in compound HO--CHO are $y$ then calculate the value of $(x+y)$?
\end{bankq}
\begin{bankq}[id=TOPS04-P2-CHEM-S1-Q11, source={TOPS - Test 04 (18 Dec 2022)}, paper=2, qno=11, type=single, subtopic={Nomenclature rules}, answer={C}, answer_src=key, difficulty=2, diff_parts={K1 S1 I0 C1 F0}, diff_note={concepts: alkyl groups | steps: isomeric groups count; check each}]
Select the correct statement -
\begin{qoptions}
\opt{A}{Homologoues can not be represented by same general formula}
\opt{B}{$\mathrm{C_5H_{11}}$--- has seven alkyl groups}
\opt{C}{$\mathrm{C_4H_9}$--- has four alkyl groups}
\opt{D}{\bfigi{138.6}{TOPS04_p2_chem_s1_q11_D.png}}
\end{qoptions}
\end{bankq}
\begin{bankq}[id=TOPS04-P2-CHEM-S1-Q12, source={TOPS - Test 04 (18 Dec 2022)}, paper=2, qno=12, type=single, subtopic={Nomenclature rules}, answer={A}, answer_src=key, difficulty=4, diff_parts={K1 S2 I1 C1 F0}, diff_note={concepts: E Z nomenclature | steps: priorities on each double bond; configuration; name}]
Select the correct name for the following diene
\bfig{89.9}{TOPS04_p2_chem_s1_q12_1.png}
\begin{qoptions}
\opt{A}{2-(E)--4--(Z)--octa--2,4--diene}
\opt{B}{2-(E)--4--(E)--octa--2,4--diene}
\opt{C}{2-(Z)--4--(E)--octa--2,4--diene}
\opt{D}{2-(Z)--4--(Z)--octa--2,4--diene}
\end{qoptions}
\end{bankq}

\section*{TNAS papers}
\begin{bankq}[id=TNAS06-P1-CHEM-S4-Q04, source={TNAS - Test 06 (30 Jul 2023)}, paper=1, qno=4, type=single, subtopic={Nomenclature rules}, answer={4723}, answer_src=solved, working={(a) 4 (b) 7 (c) acetylacetone gives 2 $\mathrm{CHI_3}$ (d) 3-ethyl-2,2,5,6,6-pentamethylheptane, ethyl at 3}, difficulty=8, diff_parts={K2 S3 I2 C2 F0}, diff_note={concepts: four short problems on isomers, iodoform and naming | steps: hydrogenation to the cyclobutane; isomers of the formula; iodoform count; IUPAC name; locant; combine}]
(a) Find out the total compounds of molecular formula $\mathrm{C_6H_{10}}$ which react with 1 eq. of $\mathrm{H_2}$/Pt to produce 1,2-dimethyl cyclobutane, without counting stereoisomers.\par
(b) Total number of structural isomers of $\mathrm{C_4H_{10}O}$.\par
\bfigi{58.5}{TNAS06_p1_chem_s4_q04_1.png}$\xrightarrow[\mathrm{I_2}]{\mathrm{KOH}}$ number of iodoform molecules produced per molecule.\par
(d) Position of ethyl group in correct IUPAC name of following compound.\par
\bfig{50.2}{TNAS06_p1_chem_s4_q04_2.png}
Write answer of part (a), (b), (c) \& (d) in the same order and present the four digit number as answer .\par
For example : If all these answer are 9 then fill 9999 .
\end{bankq}
\end{document}
